\slajdid{08-s012}
\begin{frame}[label=08-s012]
\gusttekst
Analizu logičkih kola mi radimo na tipičnom predstavniku familije. Međutim u proizvodnji je nemoguće postići da sva logička kola koja su šematski identična i imaju identične karakteristike. U tom slučaju analiza ide u dva koraka\par
\hspace{4mm}Odrede se ovi opsezi na tipičnom predstavniku\par
\hspace{4mm}Prouči se uticaj parametara na ove opsege i metodom najgoreg slučaja odrede inženjerske granice.\par\medskip
\begin{tabularx}{\textwidth}{@{}lX@{}}
$V_{IL}=\dfrac{V_{CC}}2-\Delta L$ & L – Low – maksimalan napon koji logičko kolo još uvek prepoznaje kao logičku nulu\\[5pt]
$V_{IH}=\dfrac{V_{CC}}2+\Delta H$ & H – High – minimalan napon koji logičko kolo još uvek prepoznaje kao logičku jedinicu
\end{tabularx}\par\medskip
Katalog – Data sheet – karakteristike daje proizvođač komponente logičkog kola\par
\centering{\setlength{\tabcolsep}{10pt}\begin{tabular}{lcccc}\toprule
 & Min & Typ & Max & Units\\\midrule
$V_{IL}$ & & & 0.8 & V\\
$V_{IH}$ & 2 & & & V\\\bottomrule
\end{tabular}
}
\par\medskip\raggedright
\alert{Uočiti da $\Delta L$ i $\Delta H$ ne moraju biti pozitivni odnosno da je uvek $V_{IH}$ veće od polovine napona napajanja a $V_{IL}$ uvek manje. Česta situacija kod logičkih kola sa bipolarnim tranzistorima je da je $V_{IH}$ manje od polovine napona napajanja.}
\end{frame}
